\chapter{Time-ordered products at a cut}\label{ch:13}
In the mixed response of Chapter~\ref{ch:12} the point where two pulses
meet is integrated against smooth solutions. A pair of quantum vertices is
different. Even if each vertex has already been defined as a Wick
polynomial, contracting two lines between the vertices asks for the square
of a singular propagator. Putting the vertices on opposite sides of a drawn
line does not remove the issue: both points can approach the same point of
the line. What does the regional calculation have to reproduce there?

\section{Causal ordering fixes the separated-vertex answer}
A time-ordered product $T_n(F_1,\ldots,F_n)$ is a symmetric multilinear
operator-valued distribution in its bosonic local insertions. More
precisely, its kernels are distributions; the displayed product denotes
their evaluation on smooth compactly supported couplings. An important
part of its definition is \emph{causal factorization}. If the supports of
a first group do not meet the causal past of the supports of the remaining
group, then
\begin{equation}
 T_n(F_1,\ldots,F_n)=
 T_k(F_1,\ldots,F_k)\,T_{n-k}(F_{k+1},\ldots,F_n).
 \label{eq:13-causal}
\end{equation}
The product on the right is the ordinary quantum product, with the later
group to the left. For spacelike-separated groups the opposite ordering is
also possible and the products commute. For two fields in a specified
state, the expectation of this ordered operator product is $W(x,y)$ when
$x$ is later, and $W(y,x)$ when $y$ is later;
at spacelike separation the two agree. It does not say to put all the
vertices on the left of a \emph{spatial} division later than all those on
the right. With lower time-ordered products known, causal factorization
constructs the next one off the total coincidence diagonal in an ordinary
globally hyperbolic bulk. Extending to the diagonal is a separate step
\cite{BF,HWT,FR}.

For instance, introduce smooth parts $\chi_1+\chi_2=\chi$ of the
switching function, supported in two overlapping neighborhoods of a cut.
Let $F_a=\lambda\int\chi_a\,:\!\phi^3\!:\,dV/3!$. The quadratic Dyson
term for $F_1+F_2$ is
\[
 -\tfrac12 T_2(F_1,F_1)-T_2(F_1,F_2)-\tfrac12 T_2(F_2,F_2).
\]
The cross term has coefficient one, not one half. Its contractions require
the two-point function \emph{between} neighborhoods. Matching the
commutator and the two separately renormalized regional Wick cubes need
not supply it. A smooth partition is a way of labeling the vertices. It
introduces neither a new boundary condition nor a discontinuous coupling;
the answer for a transparent cut must be independent of the partition when
the pieces are added back together.

\section{One loop with its normalization}
Use a four-dimensional flat interior patch to see the extension explicitly.
Take the massless scalar vacuum and use its normal ordering in this
calculation. For $z=x-y$ set $z^2=-(z^0)^2+|\mathbf z|^2$. In units
$\hbar=1$, our time-ordered contraction is
\[
 G_F(z)=\frac{1}{4\pi^2}\frac{1}{z^2+i0},\qquad
 \Box G_F=i\delta^{(4)} .
\]
The $i0$ is a boundary value, not a small real cutoff. To check the source
sign, spatial Fourier transformation gives
$e^{-i|\mathbf p|\,|t|}/(2|\mathbf p|)$, whose first time derivative has
jump $-i$ at $t=0$; $\Box$ contains $-\partial_t^2$.
Appendix~\ref{app:C} gives the slightly different boundary value for the
unordered two-point function.

Take two cubic vertices at $x,y$. Choosing two of the three fields at each
point and pairing them across the vertices gives
$\binom32^2 2!=18$ double contractions. Away from coincidence this part of
the second-order operator in \eqref{eq:12-dyson} is
\begin{equation}
 -\frac{\lambda^2}{4}
 \int dV_x\,dV_y\,\chi(x)\chi(y)\,
  G_F(x-y)^2\,:\!\phi(x)\phi(y)\!: .
 \label{eq:13-fish}
\end{equation}
Indeed the coefficient before counting pairings is
$(-i\lambda)^2/[2!(3!)^2]$. The remaining fields in
\eqref{eq:13-fish} have \emph{not} been contracted with external
detectors. This is the two-line loop, or fish. On the punctured relative
space $z\ne0$ its scalar distribution is
$t^\circ(z)=(16\pi^4)^{-1}(z^2+i0)^{-2}$. This notation includes the
noncoincident light cone, where the boundary value is singular but has a
consistent frequency direction. Under $z\mapsto a z$, $a>0$, the degree
is $-4$. The normal space to $x=y$ has dimension four, so the divergence
is logarithmic. An ordinary integral using the pointwise formula at the
origin would not define the coefficient.

Here one can write an extension rather than just count the free parameters.
Take $\ell>0$ of dimension length and set
\begin{equation}
 t_\ell(z)=-\frac{1}{64\pi^4}\Box
 \left[\frac{\log((z^2+i0)/\ell^2)}{z^2+i0}\right].
 \label{eq:13-fish-extension}
\end{equation}
Off the origin the bracket is defined by the same complex boundary value;
its scaling degree is two, so it extends uniquely there without increasing
its scaling degree. The derivative is then taken on distributions. To see
why the coefficient works, for $s=z^2$ one has
$\Box h(s)=4s h''(s)+8h'(s)$. For $h(s)=\log(s/\ell^2)/s$ this is
$-4/s^2$; the identity continues to the prescribed boundary value away
from $z=0$. Hence \eqref{eq:13-fish-extension} agrees with $t^\circ$ and
has scaling degree four. Two such power-counting extensions can differ by
$c\,\delta^{(4)}(z)$, but not by derivatives of that delta. In particular
the displayed family gives
\[
 t_{\ell'}-t_\ell=\frac{i}{8\pi^2}
       \log(\ell'/\ell)\,\delta^{(4)}(z),
\]
by $\Box(1/(z^2+i0))=4\pi^2 i\delta^{(4)}(z)$. Inserting a contact
difference $c\,\delta^{(4)}$ into \eqref{eq:13-fish} changes the full
operator by $-c\lambda^2\int\chi^2:\!\phi^2\!:\,dV/4$. Its mixed part
for the smooth partition $\chi=\chi_1+\chi_2$ is
\[
 -\frac{c\lambda^2}{2}\int_M\chi_1\chi_2:\!\phi^2\!:\,dV .
\]
This contact is exactly zero for disjoint supports, even though the
nonlocal fish contraction can still connect the two supports. The bulk
mass normalization matters for overlapping supports and for the full
coefficient, not as a new sheet of mass on the cut. Reality and unitarity
further restrict admissible constants in the chosen propagator convention.

\section{Power counting in the normal directions}
The fish displays the general geometry of power counting. Near the diagonal
of $n$ bulk vertices in spacetime dimension $D$, hold the common point
fixed and scale the $k=(n-1)D$ relative coordinates to zero. If the
relative distribution has scaling degree $\rho$, an extension preserving
that degree is unique for $\rho<k$. Otherwise the local ambiguity consists
of derivatives of a delta in the relative coordinates of orders
$|\alpha|\le\lfloor\rho-k\rfloor$; tangential coefficients are restricted
by the normalization axioms. In a scalar graph in $D>2$, $I$ singular
internal lines of leading degree $D-2$ and $b$ derivatives on such lines
give the superficial bound
\[
 \rho-k\ \le\ (D-2)I+b-D(n-1)=D L-2I+b
\]
for a connected cluster with $L=I-n+1$. This counts allowed local terms;
tensor cancellations can lower the actual degree. Proper subclusters must
first have their products defined, and their extension choices propagated
to larger clusters. On a curved smooth bulk the short-distance counting
is local, although allowed coefficients can now involve the metric, mass
and curvature \cite{HWick,HWT}.

In a geometry with distinguished faces the collision set can itself have
strata. For the stratum with all $n$ bulk points coincident \emph{on} a
face of codimension $q$, the normal dimension in $M^n$ is
$(n-1)D+q$: one must bring the common point to the face as well as bring
the points together. Vertices constrained from the start to the face have
a different configuration space, and hence a different normal dimension.
On a smooth stratum with a normal scaling-degree bound, possible additional
ambiguities are normal derivatives of deltas supported there, subject to
that bound. If the coefficients themselves are singular at an edge or
corner, that smaller stratum needs its own analysis; one cannot count a
bulk delta once and call it a boundary counterterm. This reasoning
classifies local freedoms \emph{if} the respective distributions and
extensions exist. Bulk theorems on smooth globally hyperbolic manifolds do
not supply a renormalization theorem for every proposed boundary problem.

For the flat fish just computed, write $X=(x+y)/2$. The uncut distribution
is independent of the coordinate $X_\perp$ normal to a planar cut.
Under joint scaling of $(z,X_\perp)$ its degree is still four, whereas
$\delta^{(4)}(z)\delta(X_\perp)$ has degree five. Thus the latter is not
an additional ambiguity with this scaling bound. More generally, an
\emph{artificial} wall is absent from the whole theory's local background
data. Regional subtractions, including any auxiliary sheet terms needed by
a regulator, must combine into the same whole-spacetime prescription at a
transparent seam. A surviving wall coupling that changes correlators when
the plane is moved fails this comparison. A \emph{physical} defect, by
contrast, may have a prescribed action, a jump equation or its own fields.
The face and its geometry then are local data of the theory. Face, corner
or junction operators allowed by power counting and symmetries must be
fixed by the physical normalization conditions when needed. Their
coefficients cannot be inferred by demanding that this defect disappear.

\section{Ward identities are more than matching conditions}
Time-ordered products obey differential identities with contact terms.
For the free scalar, $P_xW(x,y)=0$ but $P_xG_F(x,y)=i\delta(x,y)$.
Differentiating the time ordering has detected the equal-time commutator.
Treating time-ordered fields as classical solutions at the coincident
point would miss it. At zero coupling the flat massless scalar has a shift
symmetry $\phi\mapsto\phi+\epsilon$, with current
$j^\mu=-\partial^\mu\phi$ and $j^0=\partial_t\phi$. In the interior the
Ward identity for one field insertion is
\[
 \partial_\mu^x\langle\mathcal T j^\mu(x)\phi(y)\rangle
 =-\Box_x G_F(x-y)=-i\delta^{(4)}(x-y).
\]
The sign is that of $[j^0(t,\mathbf x),\phi(t,\mathbf y)]
=-i\delta^{(3)}(\mathbf x-\mathbf y)$. Integrating the identity over a
small spacetime neighborhood of $y$ detects that the inserted field shifts
by one. A cubic potential breaks this shift explicitly; that breaking is
not a quantum anomaly. For other symmetries the contact terms involve the
variations of the insertions, together with any physical boundary fluxes.
Matching the smooth classical fluxes alone does not test the quantum
contact terms.

Two failures of a local comparison have different diagnoses. The two routes may use different allowed local terms, such as
different mass normalizations in \eqref{eq:13-fish}. Aligning the
prescriptions can remove that discrepancy. Or a symmetry condition may
itself be obstructed. In the usual local perturbative framework, when the
Ward identities hold at lower orders, their first possible local violation
$\mathcal A$ obeys a consistency condition. For a gauge theory written
with a nilpotent symmetry differential $s$ this reads $s\mathcal A=0$
modulo the relevant total derivatives. Adding an admissible local
counterterm $B$ shifts the violation by $sB$ in that quotient. At an actual
boundary a total derivative has a boundary term; it is the combined bulk
and boundary identity that must be tested. A nontrivial class, not
removable by any allowed $B$, is a quantum anomaly of that Ward identity,
not an arbitrary failure to match two choices of subtraction. Conversely,
absence of such a class does not furnish a state, a transmission Green
operator or equality of every pair of renormalized products. These are
separate inputs to an order-by-order comparison \cite{FR}.

\section{Fixed sectors and an optional weighted sum}
Suppose the classical gluing problem allows several inequivalent global
sectors: for instance, transport around a loop through the seam may depend
on transition data which cannot be removed by a gauge transformation.
Choose one such sector $s_0$, its exterior conditions and its fluctuation
problem. For insertions in the common observable domain let $\iota_{s_0}$
be the comparison from the matched regional algebra to the whole algebra,
including transport of any crossing observable. The \emph{strong,
fixed-sector} assertion at order $n$ is an equality there, for every
admissible collection of insertions,
\begin{equation}
 \iota_{s_0}\!\left(T^{\rm sewn}_{n,s_0}(F_1,\ldots,F_n)\right)
 =T^M_{n,s_0}(\iota_{s_0}F_1,\ldots,\iota_{s_0}F_n).
 \label{eq:13-fixed}
\end{equation}
Here $F_i$ are matched labels for local insertions on which both products
are defined. For numerical correlators one uses corresponding states as
well. The preceding fish calculation identifies one local ambiguity that
must be fixed for this assertion even at $n=2$; it is not a proof of
\eqref{eq:13-fixed} at all orders. In a fixed flux or fixed bundle sector
there is no instruction in causal factorization to sum over others.

A different observable may deliberately leave sector data unobserved.
First map the sector answers to a common target, for example complex
amplitudes $A_s$ with the same fixed external data. If a \emph{chosen}
measure is a finite groupoid counting measure with weights $w_s$, the
combined amplitude takes the form
\[
 A_{\rm sum}=\sum_{[s]}\frac{w_s}{|\operatorname{Aut}(s)|}\,A_s .
\]
The meaning of the denominator can be checked without a path integral. If
a finite gauge group $G$ acts on a set of presentations $X$, then dividing
$\sum_{x\in X}w_x A_x$ by $|G|$ counts an orbit $Gx$ with coefficient
$|Gx|/|G|=1/|G_x|$ for invariant $w_x A_x$. Here $G_x$ is the stabilizer,
namely the automorphism group of that presentation. If the weights are
already a measure on equivalence classes, one does not divide a second
time. For continuous labels a measure and treatment of residual symmetry
volumes must be supplied; there is no default convergent sum. A statistical
mixture of states would instead require nonnegative normalized weights.
Summing complex amplitudes and mixing states are different experiments.

For a numerical example, measure the holonomy $H=ba$ in the two-strip
circle of Chapter~\ref{ch:02}. Each of the two classes $H=+1,-1$ has a
diagonal $\mathbb Z_2$ stabilizer, so unit groupoid weights give each
class counting mass $1/2$. The normalized mean of the measured sign is
$[(+1)/2+(-1)/2]/[1/2+1/2]=0$. In a fixed periodic sector the same
measurement gives $+1$, and in a fixed antiperiodic sector it gives $-1$.
Zero is the mean for the chosen ensemble, not a third possible holonomy.

Equality in each fixed sector implies equality of common weighted sums
when those sums are defined with the same measure. Equality only after a
sum is weaker, since sector discrepancies can cancel. In particular,
releasing a common classical trace to obtain transmission, contracting
quantum insertions across a cut, and summing over sectors held fixed in
\eqref{eq:13-fixed} are three different operations.

\begin{exercise}
Suppose an order-two distribution in flat four-dimensional spacetime is a
single relative derivative $\partial_{z^\mu}t_\ell(z)$ of the renormalized
fish, constant in $X$. At the coincident-on-plane stratum, what sheet
contact term is no longer ruled out by the scaling-degree count alone?
Why does that count not allow you to introduce a new freely adjustable
sheet coupling in the transparent theory?
\end{exercise}
\noindent\emph{Hint.} Scale $z$ and $X_\perp$ together. A derivative can
raise the degree by one. Then ask what additional data a coefficient with
index $\mu$ would need, and whether differentiation of the already fixed
whole distribution leaves such a choice.
